\uuid{ncjJ}
\exo7id{4458}
\titre{exo7 4458}
\auteur{quercia}
\organisation{exo7}
\datecreate{2010-03-14}
\isIndication{false}
\isCorrection{true}
\chapitre{Série numérique}
\sousChapitre{Autre}
\cpge{mpsi-mp2i.procedes-sommatoires-discrets.a, pcsi.series-numeriques.a, ptsi.series-numeriques.a, mpsi-mp2i.procedes-sommatoires-discrets.b, pcsi.series-numeriques.b, ptsi.series-numeriques.b, bcpst2.series-reelles}
\module{Analyse}
\niveau{L1}
\difficulte{}

\contenu{
\question{
Soit $(u_n)_{n\ge 1}$ une suite positive.
On pose $v_n = \frac{1}{n(n+1)}\sum_{k=1}^n ku_k$.
Montrer que les séries $\sum u_n$ et $\sum v_n$ ont même nature et
éventuellement même somme.
}
\reponse{
$\sum_{k=1}^n v_k + nv_n = \sum_{k=1}^n u_k$.\par
         Si $\sum u_n$ converge, $\sum v_n$ converge aussi (SP majorées) et
         $nv_n \to \ell  \Rightarrow  \ell = 0$.\par
         Si $\sum u_n$ diverge et $\sum v_n$ converge, alors $nv_n \to +\infty$,
         contradiction.
}
}
